\section{Introduction}
\label{sec:intro}

A fast-weight state written from context is a $d\times d$ matrix that can,
in principle, be composed algebraically at read time instead of being
rolled out sequentially. Recent work already establishes that
matrix-valued recurrent state can escape the $\mathrm{TC}^0$ ceiling that
bounds transformers and diagonal state-space models: negative-eigenvalue
linear RNNs and RWKV-7's generalized delta rule both provably solve
NC$^1$-hard state-tracking problems (the $S_5$ word problem), each via an
$O(h)$ sequential recurrent rollout over $h$ tokens.
\citep{grazzi2025unlocking,peng2025rwkv7}
% evidence: research/ncr_separation_grounding.md Part 3(b), items 5+7
So ``a matrix-valued state can state-track'' is not new, and we do not
lead with it. What we test is a narrower, algorithmic claim: given an
already-written matrix of relation operators, can a reader recover a
depth-$h$ composition in $O(\log h)$ sequential matrix multiplications via
repeated squaring, independent of how long the write took -- and does
that reader stay \emph{exact} at depths where every sequential alternative
has already drifted into noise?

Native Composition Reads (NCR) is such a reader. A fast-weight state
$Z \in \mathbb{R}^{d\times d}$ is written in-context from relation tokens
over $K$ entities forming a single Hamiltonian cycle; the model reads
$o = \mathrm{read}(Z^h, q)$, computing $Z^h$ by binary exponentiation. A
literature sweep run for this paper (Part 2 of
\texttt{research/ncr\_separation\_grounding.md}) found no published work
combining in-context-written operators, exact (not approximate)
composition, and an $O(\log h)$ query-time read; the closest neighbors --
Fast Weight Memory's recursive, approximate, fixed-depth reads
\citep{schlag2021fwm}, Log-Linear Attention's hierarchical-summarization
compute complexity, and HOLA's exact-but-verbatim key-value cache -- each
share one axis with NCR and lack another.
% evidence: research/ncr_separation_grounding.md Part 2

We report two results, in the honest order they occurred. First, NCR
survives a fully-taught adversarial construction: at $K{=}8$, holding
in-context-written operators exact through $h^\ast{=}61$ (median recovery
$1.0$) while the strongest $O(h)$ baseline has already fallen to $0.158$
(Section~\ref{sec:win}). Second, scaling the identical mechanism up the
entity-count axis $K$ hits a wall: every tested cell at $K{=}32$ --
sixteen dimension/budget combinations plus eight further budget-rescue
cells, twenty-four in total -- fails to reach robust convergence, and
far-depth recovery reads exactly zero in every one
(Section~\ref{sec:wall}). Rather than report the wall as a ceiling, we
diagnose it: the training objective is spectrum-blind by construction
(a shallow-depth cosine loss constrains direction only), the trained
write drifts non-normal and ill-conditioned, and a post-hoc orthogonal
projection of the same learned operator -- with no retraining -- already
recovers most of the lost depth (Section~\ref{sec:diagnosis}). We
pre-registered the fix experiment's verdict bands before it ran and
report its landed verdict: FAIL -- the constrained write did not train --
conditional on an in-progress code re-audit
(Section~\ref{sec:pending}).

\textbf{Contributions.} (1) A capability separation at $K{=}8$ under a
readout and task construction closed against the two standard escape
routes for this claim shape -- argmax-decoded capacity inflation and
modular collapse of held-out depth on a single cycle -- both discussed in
Section~\ref{sec:background}. (2) An honestly reported trainability wall
at $K{=}32$, replicated across three separate probes (a dimension-ratio
grid, a mechanism analysis, and a budget-rescue dissociation probe) that
jointly close the entity-count axis under a pre-registered stopping rule.
(3) A measured, not asserted, mechanism for the wall: spectral
non-normality induced by a loss that never sees conditioning. (4) A
methods contribution, the single-Hamiltonian-cycle-plus-residue-guard
construction that keeps a depth-generalization claim from silently
collapsing modulo the cycle length. (5) A pre-registered, band-scored
report of the fix experiment's landed verdict -- FAIL (the constraint did
not train through), which the pre-registered bands distinguish from a
NULL (a write that trains but buys no depth) -- reported plainly as a
finding, not hidden.
